\documentclass[10pt]{article}
\usepackage[margin=20mm]{geometry}
\usepackage{amsmath,amssymb,mathtools,fontspec,unicode-math,xurl,hyperref}
\IfFontExistsTF{DejaVu Serif}{\setmainfont{DejaVu Serif}}{\setmainfont{Latin Modern Roman}}
\setmathfont{latinmodern-math.otf}
\setlength{\parindent}{0pt}\setlength{\parskip}{6pt}\sloppy
\title{Technical Review of
Uniqueness and Minimality of Euler’s Elastica
with Monotone Curvature}
\author{AI-assisted technical review}
\date{}
\begin{document}
\maketitle
Abstract: The reviewed source is arXiv:2402.12771v2, the author-declared accepted revision of 28 March 2025. The central minimisation arguments were checked through their free-boundary reductions, scaling, elliptic parameter maps, and length-penalised dependencies. A concrete pinned critical curve is omitted by the family displayed in Remark 4.3; this is a local Minor classification defect. Two other inconsistent displays are already resolved by correct formulas in the source’s proofs and are documented as normalization clarifications. No substantive defect in the checked central uniqueness conclusions was demonstrated.\par
\section*{Identity, scope, and conventions}
The source is T. Miura and G. Wheeler, Uniqueness and minimality of Euler’s elastica with monotone curvature, arXiv:2402.12771v2 [MW], associated with DOI `10.4171/JEMS/1708`. The arXiv record calls this the accepted version; equivalence to the unacquired publisher PDF is not established. Its arXiv revision date is 28 March 2025, while its first page separately prints March 31, 2025. Both labels are retained. Page numbers below are one-based PDF pages, equal to printed pages. The initial review read all 21 pages as text; the present continuation visually inspected source pages 1–19.\par
Detailed checks covered the first variation and no-flux condition, the reductions in Sections 3–4, the explicit elliptic maps, the half-line normalization, the pinned classification in Remark 4.3, and the strict length-monotonicity argument in Section 5. The cited bootstrap and compactness theorems and the external singular-limit results used in Section 5 were not completely reproved. The abstract record for the cited phase-transition paper was located, but attempts to obtain its full text through the web reader failed; its theorem-level verification is a remaining limit. No author code was executed. Bounded acquisition and follow-up searches did not locate a separate supplement, paper-specific implementation, or correction; this is not proof of their nonexistence. Absence of code is not a mathematical defect.\par
\section*{Checked assumptions and central proof chain}
The energy is \(B[\gamma]=\int k^2\,ds\) on arc-length \(W^{2,2}\) immersed planar curves. Clamped competitors fix both endpoints and both unit tangents; fixed-length and variable-length classes are distinct. The multiplier equation is \[2k''+k^3-\lambda k=0,\qquad E^\lambda=B+\lambda L.\] Theorem 1.1 assumes nonconstant monotone signed curvature. Theorem 1.2 additionally excludes subcritical wavelike curves. Theorem 1.4 fixes nonzero angles strictly between \(-\pi\) and \(\pi\) and only asserts uniqueness sufficiently close to the straightening limit.\par
For the free-boundary classes, endpoint tangents are horizontal while endpoint vertical positions are free. In the boundary first variation, a test \(\eta=\varphi N\) with \(\varphi(0)=1\), \(\varphi'(0)=0\) isolates \(2k'(0)\), giving the no-flux condition; the other endpoint is similar. Classification then confines nonconstant minimisers to complete curvature-monotonicity intervals. The fixed-length scaling rule is \(B[\Lambda\gamma(\cdot/\Lambda)]=\Lambda^{-1}B[\gamma]\). The repeated-half-period reduction strictly reduces energy when the number of repetitions exceeds one, while preserving horizontal extent and endpoint tangents. Fixing both endpoint positions removes the residual vertical translation/reflection freedom. Replacement of a clamped subarc preserves \(W^{2,2}\) admissibility; equality competitors would contradict uniqueness of the full arc.\par
The orbitlike parameter map \[f_o(m)=\frac{2E(m)}{mK(m)}+1-\frac2m\] is strictly decreasing: with \(g=(1-m)K^2-E^2\), one has \(g(0)=0\), \(g'=-(K-E)^2/m<0\), and \(f_o'=g/[m^2(1-m)K^2]<0\). For its length-penalised counterpart the source’s numerator comparison is supported by the exact identity \[(m^2-8m+8)^2(1-m)-(4m^2-12m+8)^2=m^4(1-m)>0.\] These computations use \(0<m<1\). The wavelike penalised map is decreasing on \([m_0,1)\) because \(2E-K\leq0\) is decreasing while \(\sqrt{4m-2}>0\) is increasing. Thus the restriction excluding subcritical waves is material.\par
For Lemma 5.3, the two minimality inequalities first imply monotonicity of length with \(\varepsilon\). Equality of lengths forces equality of bending energies and makes one nonstraight curve critical for two distinct multipliers. Subtraction of the multiplier equations gives \((\lambda_1-\lambda_0)k=0\), a contradiction. This verifies the strictness needed for the inverse length map. The boundary-layer convergence and generic-angle assertions remain conditional on the cited singular-limit results, whose full proofs were not independently checked.\par
\section*{A local pinned-classification defect}
Equation reference: bad. Problem 1 (Remark 4.3 omits an actual pinned critical curve).  On pp. 16–17 the source asserts that every nonstraight pinned critical point of \(E^\lambda\) is an isometric image of \[
 \gamma(s)=N\sqrt{\frac{4m-2}{\lambda}}\,
 \gamma_w\!\left(
 N^{-1}\sqrt{\frac{\lambda}{4m-2}}\,(s-K(m)),m\right),
 \quad 0\leq s\leq2NK(m),\] where \(m\in(1/2,1)\), \(N\in\mathbb N\), and \(|2(2E-K)\sqrt{4m-2}|=\ell\sqrt\lambda/N\). The family as displayed does not classify the stated critical points.\par
\subsection*{A fully admissible counterexample}
Equation reference: true. Set \[m=\frac{51}{100},\quad \lambda=1,\quad N=2,\quad a=\frac15,\quad
 K=K(m),\quad E=E(m),\] and use the source’s exact wavelike parametrisation from Theorem 2.4: \[
 \Gamma(s)=a\bigl(\gamma_w(s/a-K,m)-\gamma_w(-K,m)\bigr),
 \qquad 0\leq s\leq 4aK.\] It is smooth and has unit speed. Since \(\operatorname{cn}(-K,m)
=\operatorname{cn}(3K,m)=0\), its endpoints are \((0,0)\) and \((\ell,0)\), where \[\ell=4a(2E-K)>0.\] Positivity needs no floating-point elliptic evaluation. From \(\sqrt{1-t}\geq1-t\) for \(0\leq t\leq1\) and the integral definitions, \[E\geq\frac\pi2(1-m/2),\qquad
 K\leq\frac{\pi}{2\sqrt{1-m}}=\frac{5\pi}{7}.\] Consequently \[2E-K\geq \left(\frac{149}{200}-\frac57\right)\pi
 =\frac{43\pi}{1400}>0.\] The curvature is \(k(s)=2\sqrt m\,a^{-1}\operatorname{cn}(s/a-K,m)\). It vanishes at both endpoints and satisfies the multiplier equation with \((4m-2)/a^2=1\). For pinned variations the endpoint displacement is zero; the remaining boundary term is \([2kN\cdot\eta']_0^L=0\). Thus [true] (\#true) is an actual critical point for all admissible pinned variations, with exactly the positive multiplier and nonnegative endpoint distance required in the remark.\par
\subsection*{Why the displayed family cannot contain it}
Every member of [bad] (\#bad) has unit speed, so its geometric length is \(2NK(m)>\pi\), independently of isometry. In contrast, \[L[\Gamma]=\frac45K(51/100)\leq\frac{4\pi}{7}<\pi.\] This proves omission even if arbitrary reparametrisation is allowed. There is also an independent multiplier mismatch: the scale \(N\sqrt{(4m-2)/\lambda}\) in [bad] (\#bad) gives multiplier \(\lambda/N^2\), rather than \(\lambda\) when \(N>1\). The executed rational check records \(1\) for the multiplier of [true] (\#true) and \(1/4\) for the displayed \(N=2\) scale.\par
\subsection*{Repair boundary and impact}
The natural repeated-wave construction instead uses \[a=\sqrt{\frac{4m-2}{\lambda}},\qquad
 a\bigl(\gamma_w(s/a-K,m)-\gamma_w(-K,m)\bigr),
 \quad 0\leq s\leq2NaK.\] Here \(N\) changes the number of half-periods, not the multiplier. Allowing a horizontal reflection, its endpoint-distance relation agrees with the relation already printed in the source: \[\left|2(2E-K)\sqrt{4m-2}\right|=\frac{\ell\sqrt\lambda}{N}.\] The absolute-value bars are present in the original page image. The earlier text-extraction-based report incorrectly described them as missing; that allegation is withdrawn. The correction established here concerns the curve’s scaling, phase and parameter interval, while retaining the published endpoint-distance equation. This does not establish every root-count or zero-distance case in the remark.\par
The discrepancy is Minor: the false assertion is an additional classification of pinned critical points. The proofs of Theorems 1.1, 1.2, and 1.4 do not use this remark, and the fixed-multiplier formulas actually used in Theorems 4.1 and 4.4 have the correct scaling. The issue is not downgraded because a correction is short; its demonstrated dependency is local. No journal-version persistence or author adoption is asserted.\par
\section*{Normalization clarifications already resolved in the proofs}
These are not counted as additional substantive findings. On p. 9, Theorem 3.1 describes the inverse of \(f_w(m)=2E(m)/K(m)-1\) as increasing and reverses the signs of \(\ell\) relative to \(m-m_0\). In fact \(E'(m)<0<K'(m)\), so \(f_w'<0\); equation (3.3), p. 11, and the actual supercritical argument on p. 16 use the correct relation. The resulting unique-parameter argument requires injectivity, which remains valid.\par
The display in Theorem 3.7, pp. 13–14, omits the \(\lambda\)-dependent scale, whereas its proof on p. 15 explicitly supplies \(a=\sqrt{2/\lambda}\) and \(aM\gamma_b(s/a)\). A direct global check confirms that proof. For a finite-energy curve, \(T+e_1\in H^1(0,\infty)\), hence \(T\to-e_1\). A continuous angle lift starts at \(0\) and tends to an odd multiple of \(\pi\). Therefore \[\widetilde E^\lambda
 =\int_0^\infty\bigl(\theta'^2+\lambda(1+\cos\theta)\bigr)\,ds
 \geq 2\sqrt\lambda\int_0^\infty
       |\theta'|\sqrt{1+\cos\theta}\,ds
 \geq4\sqrt{2\lambda}.\] Equality requires the nearest half-turn, monotonicity, and \(\theta'=\pm\sqrt{\lambda(1+\cos\theta)}\). Its solutions give precisely \(aM\gamma_b(s/a)\). For \(\lambda=8\), the unscaled curve has energy \(20\), while \(a=1/2\) gives \(16\). Thus the omitted normalization cannot be used literally, but the source’s own proof already establishes the correct result needed by Corollary 3.8. No damage to that corrected central argument has been demonstrated.\par
\section*{Conclusion}
Item | Judgment | Verified scope and remainder\par
Remark 4.3 | Minor | An actual pinned critical curve is omitted; scaling and interval repaired; published endpoint relation retained.\par
Central minimisation chain | No issue found in checked steps | External singular-limit and regularity proofs remain dependencies.\par
Displays in Theorems 3.1 and 3.7 | Clarification | Correct relations are already present in their proofs.\par
Independent acceptance | Pending | This is an initial-review continuation; no human review is recorded.\par
Scientific decision is pending; no human review or author adoption is recorded.\par
9 T. Miura and G. Wheeler, Uniqueness and minimality of Euler’s elastica with monotone curvature, arXiv:2402.12771v2, 28 March 2025. https://arxiv.org/abs/2402.12771v2. EMS Press, metadata for DOI 10.4171/JEMS/1708, consulted 27 October 2026. https://ems.press/journals/jems/articles/14299035. T. Miura, Elastic curves and phase transitions, Math. Ann. 376 (2020), 1629–1674; arXiv:1710.05890v2, 15 October 2018. Abstract/version record consulted; full text unavailable through the attempted reader.\par
\section*{Supporting checks}
4defe522864da93ab6813cb3028141c8451526ab749b68629bdb1290ec3e007d.\par
Both layout and raw Poppler extraction were used for the critical display. The initial executor performed the scientific reading, derivations, independent rational computation, report writing and initial compilation. Those unchanged calculations are retained with their original execution records. This continuation checked the existing source evidence visually and corrected the unsupported missing-absolute-value allegation. The existing native helper `/root/bounded\_0331\_check` supplied a bounded analytical challenge for this paper’s half-line and pinned-curve arguments; its identifier reflects the original helper assignment. It did not inspect images, run code, or perform independent acceptance. No external Chat or Work session was created for that check.\par
The controller’s reported UI observation and this runtime’s self-description are recorded separately; no backend equivalence is asserted. The earlier image-reader failures and uninspected render records remain preserved as historical evidence. Current viewing and final-build facts are recorded separately under recovery identifier `[internal task]`; no historical check is relabelled as newly performed. The corrected initial report remains subject to independent acceptance.\par
\end{document}
